\documentclass[11pt]{article}

\usepackage[margin=1in]{geometry}
\usepackage[T1]{fontenc}
\usepackage[utf8]{inputenc}
\usepackage{titlesec}
\usepackage{enumitem}
\usepackage{hyperref}
\usepackage{parskip}
\usepackage{microtype}
\usepackage{booktabs}
\usepackage{xcolor}
\usepackage{fancyhdr}
\usepackage{array}
\usepackage{tabularx}

% Header / footer
\pagestyle{fancy}
\fancyhf{}
\rhead{\small Point Taken AI Co-Player}
\lhead{\small Internal Documentation}
\cfoot{\thepage}
\renewcommand{\headrulewidth}{0.4pt}

% Section formatting
\titleformat{\section}{\large\bfseries}{\thesection.}{0.5em}{}[\titlerule]
\titleformat{\subsection}{\normalsize\bfseries}{\thesubsection}{0.5em}{}
\titlespacing{\section}{0pt}{1.4em}{0.5em}
\titlespacing{\subsection}{0pt}{1em}{0.3em}

% List spacing
\setlist[itemize]{topsep=3pt, itemsep=2pt, parsep=0pt}
\setlist[enumerate]{topsep=3pt, itemsep=2pt, parsep=0pt}

\hypersetup{
  colorlinks=true,
  linkcolor=black,
  urlcolor=blue,
  pdftitle={Point Taken AI Co-Player Documentation},
}

\newcolumntype{L}[1]{>{\raggedright\arraybackslash}p{#1}}

% ---------------------------------------------------------------
\begin{document}

\begin{center}
  {\LARGE \textbf{Point Taken AI Co-Player Documentation}} \\[0.5em]
  {\normalsize Internal Technical Reference} \\[0.2em]
  {Last Updated: \small \today}
\end{center}

\vspace{0.5em}
\hrule
\vspace{1.5em}

% ---------------------------------------------------------------
\section{Purpose}

The AI Co-Player is a feature for Point Taken that enables solo practice against an AI
opponent. The primary goal is to provide users with accessible practice opportunities when no
human co-player is available, and to offer a lower-pressure learning environment for new players
before they enter live multiplayer sessions.

The research goal is to evaluate whether structured solo practice against an AI opponent
improves player readiness and argument quality in subsequent human matches.

% ---------------------------------------------------------------
\section{Research Objective}

The core objective is to study whether AI-mediated solo practice improves debate performance
in live multiplayer sessions.

\textbf{Core hypothesis:} Players who complete solo practice sessions against the AI Co-Player
will produce higher-quality arguments and fewer reasoning errors in subsequent human matches
compared to players who receive no prior practice.

\textbf{Primary evaluation targets:}
\begin{itemize}
  \item Argument quality and reasoning error rate in the first human match after practice
  \item Session completion rate for solo practice
\end{itemize}

\textbf{Secondary evaluation targets:}
\begin{itemize}
  \item AI response realism (does the AI behave like a plausible typical player?)
  \item User perception of practice usefulness
  \item Retention of players who practice before their first live game
\end{itemize}

% ---------------------------------------------------------------
\section{Current Technical Stack}

\begin{itemize}
  \item \textbf{Frontend:} Nuxt 3 / Vue 3
  \item \textbf{Backend:} Node.js / Express / Socket.io
  \item \textbf{AI service:} FastAPI (Python)
  \item \textbf{LLM provider:} Anthropic Claude (\texttt{claude-sonnet-4-5})
  \item \textbf{Data store:} DynamoDB
\end{itemize}

\textbf{High-level request flow:}

\begin{center}
  \texttt{Frontend \(\rightarrow\) Backend (Socket.io) \(\rightarrow\) AI Service (FastAPI) \(\rightarrow\) Anthropic API \(\rightarrow\) Response \(\rightarrow\) Frontend}
\end{center}

The AI Co-Player extends this architecture with a dedicated co-player orchestration layer in the
backend and a new prompt pipeline in the AI service.

\textbf{Note:} The AI Feedback system (tile evaluation, 9-check fallacy detection, in-game
coaching) is already implemented and deployed on the \texttt{backend-ai-feedback} /
\texttt{demo-ai-feedback} branches. The AI Co-Player is a separate pipeline and will be built
on top of the same FastAPI service and Anthropic client.

% ---------------------------------------------------------------
\section{System Overview}

The AI Co-Player is organized around a turn-taking loop. After the human player places and
confirms a tile, the backend triggers the AI Co-Player service, which uses the current game
board and topic context to generate the AI's next move and return it as a valid tile.

Key design constraints:
\begin{itemize}
  \item The AI must follow the same rules as a human player (valid tile structure, turn order,
        no questions, no ad hominem language).
  \item The AI simulates a \textit{typical average player}, not an expert or a tutor.
  \item Responses must be topic-relevant and connect to the current thread (not generic restarts).
  \item The system must handle AI failures gracefully without breaking the session.
\end{itemize}

% ---------------------------------------------------------------
\section{AI Co-Player Behavior Specification}

\subsection{Baseline Player Simulation}

The AI Co-Player is designed to simulate an average player rather than an optimal debater.
The goal is realistic practice, not one-sided coaching.

Expected AI behavior:
\begin{itemize}
  \item Takes the opposing stance consistently throughout the session.
  \item Places arguments that respond to the current thread.
  \item Uses simple, direct language comparable to a typical student player.
  \item Occasionally makes arguable or imperfect claims (consistent with average human play).
  \item Does not lecture, tutor, or explicitly correct the human player.
\end{itemize}

\subsection{Follow-Up Turn Behavior}

When generating a response tile, the AI must:
\begin{itemize}
  \item Read the parent tile and all ancestor tiles in the current thread.
  \item Respond specifically to the parent tile's claim.
  \item Not repeat its own prior arguments.
  \item Stay on the debate topic established at session start.
\end{itemize}

\subsection{Stance}

The AI must maintain its assigned opposing stance for the full session.

\begin{itemize}
  \item If the human player is Plus (supporting the topic), the AI plays Minus.
  \item The AI may not switch positions mid-session.
  \item Rewrites and follow-ups must remain consistent with the assigned side.
\end{itemize}

% ---------------------------------------------------------------
\section{Functional Requirements}

\subsection{Core Gameplay}

\begin{enumerate}
  \item[\textbf{FR-01}] Users can start a solo practice session without waiting for a human partner.
  \item[\textbf{FR-02}] Users select the debate topic and their preferred stance before play begins.
  \item[\textbf{FR-03}] The AI automatically takes the opposing stance to the user's choice.
  \item[\textbf{FR-04}] The AI simulates baseline (average) player behavior, not an expert.
  \item[\textbf{FR-05}] AI responses are topic-relevant and connect to the current board thread.
  \item[\textbf{FR-06}] AI follows all tile structure and turn rules identically to human play.
  \item[\textbf{FR-07}] Full solo sessions complete without requiring a human partner at any point.
  \item[\textbf{FR-08}] Clear visual feedback is shown while the AI generates a response.
  \item[\textbf{FR-09}] The system recovers gracefully when the AI is slow, unavailable, or produces
        an invalid move.
  \item[\textbf{FR-10}] All turns, timestamps, topics, stances, and response metadata are logged.
  \item[\textbf{FR-11}] Anonymized session data can be exported for research evaluation.
\end{enumerate}

\subsection{Tile Reactions}

Players may react to their opponent's confirmed tiles using a predefined set of reactions
(e.g., thumbs up). Reactions are a lightweight social signal that acknowledges a strong or
interesting opposing argument without interrupting the turn flow.

\begin{enumerate}
  \item[\textbf{FR-12}] A player can react to any confirmed tile placed by their opponent.
  \item[\textbf{FR-13}] The supported reaction set for v1 is a single positive reaction
        (thumbs up / \(\mathbf{+1}\)). Additional reaction types may be introduced in later
        versions.
  \item[\textbf{FR-14}] A player may react to a given tile at most once; toggling removes
        the reaction.
  \item[\textbf{FR-15}] Reaction counts are visible to both players in real time via
        Socket.io broadcast.
  \item[\textbf{FR-16}] Reactions do not affect turn order, scoring, or game outcome in v1.
  \item[\textbf{FR-17}] In AI Co-Player sessions, the AI does not send reactions. The human
        player may still react to AI tiles.
  \item[\textbf{FR-18}] All reaction events (tile ID, actor, reaction type, timestamp) are
        logged and included in session exports.
\end{enumerate}

\subsection{Session End and Tile Limits}

Sessions do not have a fixed end condition but must not run indefinitely. The design uses a
soft pressure model: play is uninterrupted below a tile threshold; above it, both players are
increasingly nudged toward ending, with an automatic close at the hard cap.

\begin{enumerate}
  \item[\textbf{FR-19}] A configurable \textbf{soft tile limit} (default: 20 confirmed tiles)
        triggers end-game prompts. Below this threshold, no end-game UI is shown.
  \item[\textbf{FR-20}] Once the soft limit is reached, a persistent but non-blocking banner
        is shown to both players after each subsequent turn, indicating that the board is
        getting full and either player may call for a vote to end.
  \item[\textbf{FR-21}] Either player may initiate an \textbf{end-game vote} at any point
        after the soft limit is reached. The other player must accept or decline within a
        configurable timeout (default: 60 seconds).
  \item[\textbf{FR-22}] If the receiving player accepts, the session ends immediately and the
        end-state screen is shown to both players.
  \item[\textbf{FR-23}] If the receiving player declines or does not respond within the
        timeout, play continues. The requesting player may re-initiate a vote on any subsequent
        turn.
  \item[\textbf{FR-24}] A configurable \textbf{hard tile cap} (default: 30 confirmed tiles)
        triggers automatic session end regardless of player agreement. Both players are shown
        a brief notification before the end-state screen.
  \item[\textbf{FR-25}] The soft limit, hard cap, and vote timeout are configurable per
        deployment environment (not hardcoded).
  \item[\textbf{FR-26}] In AI Co-Player sessions, the AI always accepts an end-game vote.
        The AI never initiates one.
  \item[\textbf{FR-27}] Session end events (who initiated, accepted/auto, tile count at close,
        timestamp) are logged for research analysis.
\end{enumerate}

\subsection{User Experience Requirements}

\begin{enumerate}
  \item[\textbf{UR-01}] Users can practice even when no human co-player is available.
  \item[\textbf{UR-02}] New users can learn game mechanics in a lower-pressure solo environment.
  \item[\textbf{UR-03}] AI responses feel on-topic and connected to the current thread.
  \item[\textbf{UR-04}] Users are not left waiting an unreasonable amount of time for AI turns.
  \item[\textbf{UR-05}] It is always clear when it is the user's turn versus the AI's turn.
  \item[\textbf{UR-06}] The AI is perceived as following the same visible rules as human players.
  \item[\textbf{UR-07}] Users finish sessions with a clear next step (e.g., join a live game).
  \item[\textbf{UR-08}] Solo practice prepares users for subsequent human matches.
  \item[\textbf{UR-09}] The AI behaves like a plausible typical player, not an obvious AI tutor.
  \item[\textbf{UR-10}] Reacting to a tile is a single tap and does not disrupt the player's
        turn flow.
  \item[\textbf{UR-11}] Players can always see the current confirmed tile count and how far
        they are from the soft limit and hard cap.
  \item[\textbf{UR-12}] The path to ending a session is explicit and requires deliberate action;
        sessions cannot be ended accidentally.
  \item[\textbf{UR-13}] Once the soft limit is reached, the prompt to consider ending feels
        like a natural nudge rather than a disruption or failure state.
\end{enumerate}

% ---------------------------------------------------------------
\section{Technical Requirements by Layer}

\subsection{Backend}

\begin{itemize}
  \item Practice sessions are stored separately from or clearly distinguished from multiplayer games.
  \item The AI turn orchestration service receives topic, opposing stance, and full board context.
  \item State validation enforces legal turn order and valid tile structure before AI moves.
  \item All session data is logged at sufficient granularity for research analysis.
  \item Admin and export tooling provides approved access to anonymized session data.
  \item Failure recovery includes retries and safe fallback messaging for the frontend.
  \item Reaction events are stored per tile in the game record and broadcast to all session
        participants via Socket.io.
  \item Tile count is maintained server-side; soft limit and hard cap thresholds are read from
        environment configuration, not hardcoded.
  \item End-game vote state (pending, accepted, declined, expired) is tracked in the game record.
        The hard-cap auto-end is enforced server-side on every tile confirmation, not client-side.
\end{itemize}

\subsection{AI Service}

\begin{itemize}
  \item The AI prompt encodes Point Taken board structure, tile relationships, and turn conventions.
  \item Baseline player simulation is explicitly specified (average behavior, not expert).
  \item Topic relevance is maintained across the full session using accumulated context.
  \item Opposing stance fidelity is enforced at the prompt level and verified in output.
  \item Follow-up quality is measured against criteria: relevance, non-repetition, coherence.
  \item An evaluation harness supports stance fidelity, topic relevance, and realism scoring.
  \item All agent outputs must be valid JSON only; no extra text outside the JSON.
  \item The AI service uses the Anthropic Python SDK (\texttt{anthropic}) with
        \texttt{claude-sonnet-4-5} as the default model. The model is configurable via the
        \texttt{ANTHROPIC\_MODEL} environment variable.
  \item The existing AI Feedback pipeline (\texttt{agents.py}) provides reference patterns for
        prompt structure, JSON output enforcement, and fallback handling that the co-player
        pipeline should follow.
  \item The AI always accepts end-game votes; this is handled in the backend orchestration
        layer, not in the AI service prompt.
\end{itemize}

\subsection{Frontend}

\begin{itemize}
  \item A clear entry point separates practice mode from multiplayer game flow.
  \item A topic and side selection flow precedes session start.
  \item Turn state is always visible (human turn / AI thinking / AI turn played).
  \item AI loading and thinking states are visually distinct.
  \item Session end state provides clear onward options for the user.
  \item Error states display graceful fallback messaging.
  \item Each opponent tile renders a reaction affordance (thumbs up icon + count). The icon
        is visually active when the current player has reacted to that tile.
  \item A tile counter (e.g., \textit{``18 / 20 tiles''}) is visible on the board at all times.
        The counter changes color or style as it approaches the soft limit.
  \item Above the soft limit, a persistent banner prompts players to consider ending. A single
        action opens the end-game vote flow.
  \item The end-game vote UI shows clearly who initiated, and presents the receiving player
        with an accept / decline choice and a countdown to timeout.
  \item The hard-cap auto-end displays a brief ``Board full'' notification before transitioning
        to the end-state screen.
\end{itemize}

% ---------------------------------------------------------------
\section{Prompt Engineering Requirements}

Prompts for the AI Co-Player must satisfy the following:

\begin{itemize}
  \item Provide explicit game rules (tile types, turn order, stance constraints).
  \item Specify the target player persona (typical average player, not expert).
  \item Include positive examples of good AI moves and negative examples to avoid.
  \item Add hard rules to prevent off-topic responses, stance switching, and rule violations.
  \item Require structured JSON output with no additional explanatory text.
  \item Version all prompts so changes can be tracked across evaluation rounds.
\end{itemize}

The general philosophy is \textbf{realism first}. The AI should feel like a plausible debate
partner, not a helpful assistant or a content moderator.

\textbf{Reference implementation:} The AI Feedback system's \texttt{analysis\_agent} prompt
(in \texttt{agent/agents.py}) demonstrates the required output discipline: strict JSON schema,
decision-procedure format, and version-tagged prompts. Co-player prompts should follow the
same conventions.

% ---------------------------------------------------------------
\section{Evaluation Framework}

\subsection{Product Readiness Criteria}

\begin{itemize}
  \item AI responses are consistently topic-relevant across all tested debate topics.
  \item The AI maintains opposing stance fidelity throughout full sessions.
  \item AI turns feel coherent and connected to the parent tile, not generic.
  \item Session failure rate (invalid move, crash, timeout) is below an acceptable threshold.
  \item Users report the AI as a believable practice partner in usability testing.
\end{itemize}

\subsection{Operational Targets}

\begin{itemize}
  \item $\geq$85\% of AI responses rated as topic-relevant on a labeled evaluation set.
  \item $\geq$95\% opposing-stance fidelity across reviewed session samples.
  \item $\leq$10\% invalid-move rate (malformed tiles, rule violations) in stress testing.
  \item 100\% metadata presence: tileId, topic, stance, modelName, promptVersion, timestamp.
  \item AI turn latency within acceptable bounds for live session flow. ($\leq$10 seconds per response)
\end{itemize}

\subsection{Evaluation Harness}

The evaluation harness should include:
\begin{enumerate}
  \item A labeled test set of 15--25 game scenarios (topic, board state, expected AI move type).
  \item Automated checks for stance fidelity and JSON validity.
  \item Human review rubric for topic relevance and typical-player realism.
  \item Before/after outputs for each prompt iteration.
  \item Summary report of false positives (off-topic, wrong-stance, rule-violating moves).
\end{enumerate}

% ---------------------------------------------------------------
\section{Data Collection and Analysis}

The system must log sufficient data to support both debugging and research evaluation.

\textbf{Required log fields per session:}
\begin{itemize}
  \item Session ID and timestamp
  \item Debate topic and assigned stances (human and AI)
  \item Full tile sequence with authorship (human vs.\ AI) and timestamps
  \item AI model name and prompt version for each AI turn
  \item Any failure events (timeout, invalid move, fallback triggered)
  \item All reaction events: tile ID, reactor, reaction type, timestamp
  \item Session end event: initiator (human or auto), tile count at close, outcome
        (agreed / declined / hard-cap), timestamp
\end{itemize}

\textbf{Research analysis targets:}
\begin{itemize}
  \item Comparison of argument quality between players who practiced and those who did not.
  \item AI response quality over time as prompts are iterated.
  \item Session completion rate and drop-off points.
  \item Reaction rate as a proxy signal for argument quality and engagement
        (tiles that receive reactions may indicate stronger or more surprising arguments).
  \item Distribution of session lengths and end-game trigger types (agreed, declined, hard-cap)
        to inform default tile limit configuration.
\end{itemize}

% ---------------------------------------------------------------
\section{Experimental Design}

The study should support at least two player conditions:

\begin{itemize}
  \item \textbf{Control:} Players enter their first live game with no prior solo practice.
  \item \textbf{Treatment:} Players complete at least one AI Co-Player session before live play.
\end{itemize}

\textbf{Core metrics:}
\begin{itemize}
  \item Reasoning error rate in first live game
  \item Argument coherence scores
  \item User-reported confidence and preparation
\end{itemize}

\textbf{Implementation requirement:} The system must include a reliable condition assignment
mechanism so players can be consistently allocated to practice vs.\ no-practice conditions.

% ---------------------------------------------------------------
\section{Open Questions}

The following questions require alignment before full implementation:

\begin{enumerate}
  \item Should practice mode mirror the live multiplayer board exactly or use a simplified
        solo-specific flow?
  \item Is post-round review (showing AI move quality, flagging errors) part of v1 scope?
  \item What is the minimum session length that constitutes meaningful practice?
  \item Should future versions support differentiated AI player types (e.g., aggressive,
        conciliatory, evidence-focused)?
  \item How should practice sessions be surfaced to users --- opt-in, prompted, or required
        for new accounts?
  \item Should reactions be visible to the opponent in real time, or only shown as a
        summary at session end? Real-time visibility may influence player behavior (positive
        feedback loop) in ways that affect research validity.
  \item Should the reaction set expand beyond thumbs up in v1 (e.g., a ``strong point''
        badge, a ``disagree'' reaction)? Additional reaction types add signal richness but
        increase UI complexity.
  \item What are the right defaults for the soft limit (20 tiles) and hard cap (30 tiles)?
        These should be validated against pilot session data before the research study begins.
  \item Should players be able to agree to end a session \textit{before} the soft limit is
        reached (e.g., if one player wants to quit early), and if so, how is that handled
        in the research dataset?
\end{enumerate}

% ---------------------------------------------------------------
\section{Role Specifications}

\subsection{Technical Lead \& Core Developer --- Corey Zhang}

\begin{itemize}
  \item Own and implement the full AI Co-Player pipeline end-to-end: session creation,
        LLM integration, response validation, turn orchestration, and logging.
  \item Build the co-player agent in \texttt{agent/agents.py} using the Anthropic Python SDK,
        following the patterns established by the existing AI Feedback pipeline.
  \item Define and enforce the AI input/output contract across the backend and AI service.
  \item Implement the \texttt{/co-player/turn} FastAPI endpoint: receive board context, call
        the LLM, validate output, and return a well-formed tile.
  \item Build output validation and fallback handling (malformed JSON, invalid tile structure,
        stance drift).
  \item Wire the co-player endpoint into the backend Socket.io game flow in coordination
        with Christopher.
  \item Implement session logging and research export tooling; instrument the pipeline with
        prompt version, model name, latency, and failure metadata.
  \item Set prompt engineering standards: stance fidelity, topic relevance, baseline persona.
  \item Define evaluation criteria and oversee the evaluation harness.
  \item Document all prompts, test cases, and design decisions in shared project resources.
  \item Coordinate work between Kim and Mahnum; surface blockers to Steve, Athira, and Emma.
\end{itemize}

\subsection{Prompt Engineering --- Mahnum Somji}

\begin{itemize}
  \item Develop and iterate the AI Co-Player turn-generation prompt: persona specification,
        stance enforcement, topic relevance, and output schema.
  \item Iterate prompts based on evaluation harness results and false-positive reports.
  \item Maintain versioned prompt files and document each revision with before/after outputs.
  \item Build and expand the labeled evaluation harness (15--25 game scenarios minimum).
  \item Run human review passes for topic relevance and typical-player realism scoring.
  \item Deliver prompt iterations and evaluation summaries on a two-week cycle.
\end{itemize}

\subsection{Model Orchestration Support --- Kim Xiao}

\begin{itemize}
  \item Support Corey on AI service integration tasks as needed.
  \item Assist with testing and validating co-player outputs against the evaluation harness.
  \item Help instrument and review session logs for research metadata completeness.
  \item Coordinate with and mentor Mahnum on prompt testing and iteration feedback loops.
\end{itemize}

\subsection{Backend Dev --- Christopher Heo}

\begin{itemize}
  \item Implement the practice session lifecycle (creation, turn orchestration, persistence).
  \item Integrate the AI Co-Player service into the existing Socket.io game flow.
  \item Implement failure recovery: retries, timeouts, and safe fallback messaging.
  \item Build the session logging layer and admin export endpoints.
  \item Implement the tile reaction Socket.io events (\texttt{addReaction},
        \texttt{removeReaction}) and persist reaction data in the game record.
  \item Implement tile count tracking, soft-limit detection, end-game vote state machine,
        and hard-cap enforcement in the backend game flow.
\end{itemize}

\subsection{Frontend Dev + UI/UX Design --- Christopher Heo, Audrey Chung}

\begin{itemize}
  \item Build practice mode entry points and topic/side selection flow.
  \item Implement AI thinking state and turn ownership indicators.
  \item Handle error states and session end with clear next-step options.
  \item Ensure the practice board and live game board share visual consistency.
  \item Design and implement the per-tile reaction affordance on opponent tiles.
  \item Design and implement the tile counter, soft-limit banner, end-game vote UI,
        and hard-cap auto-end notification.
\end{itemize}

\end{document}
